\chapter{Conclusion and Future Work}
\section{Conclusion}
The E-ShoeS (ESS) is a smart running shoe system that integrates kinetic energy harvesting from footsteps, GPS positioning, AI-assisted emergency SOS tracking, and active lighting. Piezoelectric transducers embedded in the sole harvest mechanical energy from walking to supplement the internal 3.7 V, 4000 mAh rechargeable Li-Po battery, which provides continuous power for the lights, BLE module, and emergency alert systems.

For safety and user control, an embedded microcontroller drives the front headlights and surround RGB LED strip during night activities. When the user triggers an emergency alert, the AI feature on the smartphone app optimizes GPS updates and sends live location messages to emergency contacts. Ongoing research and development will continue to focus on improving component durability, weight distribution, and battery protection.

\section{Future Work}
Building on the analytical models and initial prototype, the next phase will focus on extended bench validation, environmental qualification, and commercial optimization:
\begin{itemize}[label=--]
    \item Power system validation: carry out extended testing of the prototype to measure actual PZT-5H piezo conversion losses, solar MPPT efficiency under dynamic foot movement and thermal dissipation through the carbon-Kevlar and aerogel battery housing under peak internal load.
    \item Durability and environmental qualification: Carry out cyclic mechanical fatigue testing on the polymer reinforced PZT harvesters (targeting> 5 million compression cycles) and the serpentine FPCB LED strips (targeting> 1 million bending cycles). This will also include qualification of the waterproof housing.
    \item Ergonomic optimisation and commercial scaling: carry out human factors gait analysis to assess user comfort and weight distribution over the course of a long running session, then a bill of materials cost optimisation to establish the Basic and Sport ranges.
\end{itemize}